% !TeX root = ../se200_referential_regimes.tex

% =============================================================================
\section{The Lower Bound}
\label{sec:bound}
% =============================================================================

The previous section collected the nine core identity regimes.
This section proves the lower-bound claim.
The proof requires two facts:
the inventory requires stable reference to six carrier kinds,
and three of those carrier kinds exhibit witnessed basis plurality under the transformation basis.

% -----------------------------------------------------------------------------
\subsection{The Inventory Floor}
\label{subsec:inventory_floor}
% -----------------------------------------------------------------------------

The reference inventory of Section~\ref{sec:inventory} requires stable reference to six kinds of carrier:
obligation-bearers, rules, occurrences, scopes, records, and plain referents.
Stable reference to a carrier kind requires an identity regime for that kind,
because the substrate must determine when representations of that carrier count as the same referent.

\begin{lemma}[Inventory floor]
\label{se200.lem.InventoryFloor}
Under the assumptions of Section~\ref{sec:formal_setting},
a neutral accountability substrate realizing the inventory of Section~\ref{sec:inventory}
must realize at least one identity regime for each required carrier kind.
\end{lemma}

\begin{proof}
By Definition~\ref{se200.def.ReferenceKind}, each reference kind in the inventory
must be identifiable, trackable, and distinguishable from other referents
without making object-level causal or normative commitments.
By Definition~\ref{se200.def.IdentityRegime}, an identity regime specifies
sameness for a referent under transformation.
Thus each required carrier kind needs identity conditions sufficient to
preserve stable reference to carriers of that kind.

The required carrier kinds have disjoint representation domains in the inventory.
A representation of an obligation-bearer is not a representation of an occurrence;
an occurrence is not a record;
a rule is not a scope;
a scope is not a plain referent;
and a record is not the thing it records.
Stable reference to each required carrier kind therefore requires identity
conditions for that carrier kind.
Collapsing carrier kinds would collapse the reference inventory.
Therefore each required carrier kind contributes at least one identity regime.
\end{proof}

The inventory therefore gives a floor of six regimes.
The witnessed basis pluralities for plain referents, scopes, and rules
raise that floor to nine.

% -----------------------------------------------------------------------------
\subsection{Witnessed Splits}
\label{subsec:witnessed_splits}
% -----------------------------------------------------------------------------

Three carrier kinds exhibit witnessed basis plurality under the transformation basis.

First, plain referents split into locus-fixed and object-fixed identity.
By Proposition~\ref{se200.prop.PlainReferentUnderdetermination},
branch/fork separates these bases.
A fork may preserve the locus while breaking object identity.
When both bases are in play,
plain referents require both $\LOC$ and $\OBJ$.

Second, scopes split into extension-fixed and structure-fixed identity.
By Proposition~\ref{se200.prop.ScopeUnderdetermination},
aggregation/decomposition separates these bases.
A decomposition may preserve the covered cases while changing the internal scope structure.
When both bases are in play,
scopes require both $\SCOPEE$ and $\SCOPES$.

Third, rules split into content-fixed and structure-fixed identity.
By Proposition~\ref{se200.prop.RuleUnderdetermination},
refinement separates these bases.
A refinement may preserve rule content while changing textual or internal structure.
When both bases are in play,
rules require both $\RULEC$ and $\RULES$.

Each witnessed split adds one regime beyond the one-regime-per-carrier floor.
The additional regimes are not new reference kinds.
They are distinct identity bases for required carrier kinds.

% -----------------------------------------------------------------------------
\subsection{The Bound}
\label{subsec:the_bound}
% -----------------------------------------------------------------------------

\begin{theorem}[Nine-regime lower bound]
\label{se200.thm.NineRegimeLowerBound}
Under the assumptions of Section~\ref{sec:formal_setting},
a neutral accountability substrate that supports the inventory of
Section~\ref{sec:inventory}
and supports the witnessed basis pluralities for plain referents,
scopes, and rules
must realize at least nine identity regimes.
\end{theorem}

\begin{proof}
By Lemma~\ref{se200.lem.InventoryFloor},
the six required carrier kinds in Section~\ref{sec:inventory}
require at least six identity regimes.

Three of those carrier kinds require an additional regime when the witnessed
basis pluralities are in play.
Plain referents require both $\LOC$ and $\OBJ$ by
Proposition~\ref{se200.prop.PlainReferentUnderdetermination}.
Scopes require both $\SCOPEE$ and $\SCOPES$ by
Proposition~\ref{se200.prop.ScopeUnderdetermination}.
Rules require both $\RULEC$ and $\RULES$ by
Proposition~\ref{se200.prop.RuleUnderdetermination}.

Therefore the required core contains:
\[
\begin{gathered}
\OBL,\quad \OCC,\quad \REC,\quad \LOC,\quad \OBJ,\\
\SCOPEE,\quad \SCOPES,\quad \RULEC,\quad \RULES .
\end{gathered}
\]
By Proposition~\ref{se200.prop.DistinctCoreProfiles},
these nine core profiles are pairwise distinct.
Thus the required regimes number at least,
\[
1 + 1 + 1 + 2 + 2 + 2 = 9.
\]
Thus any neutral accountability substrate satisfying the stated assumptions
and supporting the witnessed basis pluralities
must realize at least nine identity regimes.
This proves the lower-bound claim.
\end{proof}

% -----------------------------------------------------------------------------
\subsection{Monotonicity}
\label{subsec:monotonicity}
% -----------------------------------------------------------------------------

\begin{corollary}[Monotonicity]
\label{se200.cor.Monotonicity}
Enlarging the transformation basis cannot lower the bound
established by Theorem~\ref{se200.thm.NineRegimeLowerBound}.
\end{corollary}

\begin{proof}
The core is witnessed by distinctions
already forced by the transformation basis $\Fset$.
Adding transformations adds classification coordinates.
It may leave the existing distinctions unchanged.
It may separate an already-distinct pair in further ways.
It may also force further splits.
But it cannot erase a split already witnessed by $\Fset$.
Therefore any transformation basis containing $\Fset$ still requires
at least the nine core regimes derived here.
\end{proof}

The monotonicity claim is intentionally limited.
It establishes that incompleteness of the transformation basis
cannot make the present lower bound disappear.
A richer basis may leave the bound unchanged or raise it,
but cannot lower it.

% -----------------------------------------------------------------------------
\subsection{Boundary of the Proof}
\label{subsec:bound_not_claim}
% -----------------------------------------------------------------------------

The claim proved here is a lower bound.
Under the stated assumptions and witnessed basis pluralities,
fewer than nine regimes cannot preserve stable reference
without changing transformation behavior or reintroducing a hidden regime.

The proof does not establish an upper bound.
It does not claim that the nine regimes exhaust all coherent identity regimes.
It does not claim that every deployed system must expose these regimes.
Those broader limits are stated in Section~\ref{sec:limits}.
